\documentclass[11pt,a4paper,oneside]{book} %scrbook book report

\usepackage[
backend=biber,
style=numeric,
sorting=ynt
]{biblatex}
\addbibresource{references.bib}

% Essential packages
%\usepackage{tla}
\usepackage{amsmath, amssymb, amsthm} % AMS Packages
\usepackage{graphicx,color}           % Packages for graphics and color
\usepackage[left=1.5in, right=1in, top=1in, bottom=1in, includefoot, headheight=13.6pt]{geometry}

% Optional customization packages
\usepackage{lmodern}                  % Custom fonts
\usepackage[T1]{fontenc}              % Ensure correct font encoding
\usepackage{tlatex}
\usepackage{url}

% Customising chapter headings - sectsty.pdf
\usepackage{sectsty}
\chapterfont{\Large\sc\centering}
\chaptertitlefont{\centering}
\subsubsectionfont{\centering}

\usepackage[hang, small, bf, margin=0pt, tableposition=bottom]{caption}
\setlength{\abovecaptionskip}{10pt}   % Custom captions

% Tables
\usepackage[table]{xcolor}
\usepackage{colortbl}

\newcommand{\mc}[2]{\multicolumn{#1}{c}{#2}}
\definecolor{Gray}{gray}{0.83}

\newcolumntype{x}{>{\columncolor{Gray}}c}
\newcolumntype{y}{>{\columncolor{white}}c}

% Page layout
\usepackage{parskip}
\parindent 0pt
\parskip 1ex
\renewcommand{\baselinestretch}{1.49}
\numberwithin{equation}{section}
\renewcommand{\bibname}{References}
\renewcommand{\contentsname}{Contents}
\pagenumbering{roman}


\newcommand{\acrolabel}[1]{\makebox[3cm][l]{\textbf{#1}}}
\newenvironment{acronyms}{\begin{list}{}{\renewcommand{\makelabel}{\acrolabel}}}{\end{list}}

% \includeonly{tex/chapter1}          % Option to generate specific chapters

% Customising headers - fancyhdr.pdf
\usepackage{fancyhdr}
\pagestyle{fancy}
\rhead{}
\lhead{\nouppercase{\textsc{\leftmark}}}
\renewcommand{\headrulewidth}{0pt}
\makeatletter
\renewcommand{\chaptermark}[1]{\markboth{\textsc{\@chapapp}\ \thechapter:\ #1}{}}
\makeatother


% Hyperreferencing and citations
\usepackage{hyperref}
\hypersetup{
    colorlinks,
    citecolor=blue,
    filecolor=blue,
    linkcolor=blue,
    urlcolor=blue
}
%\usepackage{natbib}
%\usepackage{apacite}
%\bibliographystyle{apacite}

% Section symbol
\usepackage{cleveref}
\crefname{section}{\S}{\S\S}
\Crefname{section}{\S}{\S\S}
\crefname{subsection}{\S}{\S\S}
\Crefname{subsection}{\S}{\S\S}


\usepackage{tabu}
\usepackage{adjustbox}
\usepackage{booktabs}% for better rules in the table
\usepackage{graphicx}
%\usepackage{subfigure}
%\usepackage{color}
%\usepackage{colortbl}
%\usepackage{soul}
%\usepackage{listings}
%\lstloadlanguages{Java,XML}
%\lstset{frame=lines}
\usepackage{./styles/astron}
%\usepackage{xspace}
%\usepackage[leqno]{amsmath}
%\usepackage{hyperref}
%\usepackage{sfmath}
%\usepackage{setspace} 
\usepackage{./styles/nust}
\graphicspath{{./img/}{./figs/}}
\usepackage[shortlabels]{enumitem}

\title{A SPECIFICATION OF A MULTI-VERSION STORAGE ENGINE USING TLA+}
%\subtitle{A sub-title}

\author{Goodwill T. Khoa}
\regno{Student Number: 219080534}
\degree{\BCHS} % MSCSE, MSCCS, BCHS
\school{\NUST}

\adviser{Prof. J. Quenum}
\adviserAffiliation{Department of Computing and Informatics }

\date{ 07 October 2022}

% \setcounter{tocdepth}{2}
% \setstretch{1.1}
% \linespread{1.1}

\begin{document}
\maketitle

\thesisAcceptanceCertificate
%  This below  can be commented out
\evaluationcommitteeapproval{ }{ }{ }

\chapter*{Dedication}
This thesis is dedicated to all the deserving children who do not have access to quality education especially young girls.

\certificateoforiginality

\acknowledgement{To God Almighty for the strength to never give up and pursue all things towards His Kingdom and Glory. To my family who through all sleepless nights, never complained about my present-absence. My wife Lutopu and son Michael and daughter Lillian as well as newly adopted children.To my Work colleagues for patience and support, To Prof J. Quenum for patience and faith in me. Thank you all.  
}

\tableofcontents
\listoffigures
\listoftables
% \lstlistoflistings

\chapter*{Abstract}
It is well known that; within complex systems and concurrency, fault deduction and system
failures is rated as the greatest risk factor amongst software developers. Even with various
efforts to eliminate some of the expected errors and or bugs. Distributed Systems tend not to
be for faint-hearted software developers. As a result this working document will focus on the
ability to check the correctness of our system (a distributed storage engine) beyond testing.
One such solution is to use the Temporal Logic of Actions (TLA) in the designing and developmental processes before actual coding takes place.  Temporal Logic of Actions Plus (TLA+) is the specification
language used throughout this research process. The specification is focused on the key
operations of the storage engine: inserting, updating and deleting and reading data from either a distributed physical storage medium or a local physical storage medium.

\resetpagenumbering

\chapter{Introduction and Motivation}\label{c-intro}

\section{Introduction}
Currently cloud computing and outsourcing is a general norm amongst developers. African,
and Namibian software developers to be precise, finds themselves with a hot bill to swallow with
regards to outsourcing storage, either for personal projects or cliental. It is with this in mind
that a small and compact team from the Namibia University of Science and
Technology embarked on the development of a storage engine, that aims to solve some of the basic but crucial architectural aspects of the development and deployment process within our country and continent as a whole.
According to the Wikipedia contributors (2022), “A database engine (or storage engine) is the
underlying software component that a database management system (DBMS) uses to create,
read, update and delete (CRUD) data from a database.” 

%Skipping Paragraphs by one line
\setlength{\parskip}{10pt}

Storage engines can be classified as
either as transactional or non-transactional, and categorized as local or distributed. The storage engine under the microscope, is a transactional and distributed storage engine which aims to have high availability with high scalability as well as strong consistency and can be used for data analysis. A brief highlight of the level of complexity within the distributed storage engine includes; Conversions, Version Control, Synchronizations, Concurrent Read-Write operations, Transaction Support, Backups and Quality Assurance, High availability and reduced time complexity in terms of overall performance. A similar dilemma was faced by Senior Developer
at Amazon Web Service (AWS). It is with this in mind that the researcher intends to utilize TLA+ a formal specification method to specify the key operations of the multi-versioned storage engine.  According to Lamport (2003). "The language of PlusCal expressions is TLA+, a high-level specification language based on set theory and first order logic.
The researcher assumes that the reader has a background in Mathematics with specific reference to the Set theory and Set Notation as the main translation of the derived algorithm will be denoted in Set Theory Notation. Below are small snippets of what to expect through out the document.

\begin{itemize}
\item Actions: Å, M, M1, M2
\item Predicates: P, Init \baro
\item Variables: Path, VersionNumber, Pointer
\item Primed variables: rootPath` PreVersion`, PPoniter`
\item States: s, s`
\item State function: f
\item Behavior: <s0, s1,..., sn>
\item Values: Data items, e.g. Integers, constants, String, Boolean
\item Semantics: [f],[Pointer], K[VersionNumber], rootPath[f]
\item Formulas: F, G, \textPhi
\item Operators: \Box, \diamond, \neq, \vee, \wedge %done
\item Quantifiers: \forall, \exists %done
\item Symbols:  \triangleq, <, >, \backsim, [, ], (, ), =, \equiv, \cdots

\end{itemize}

% Problem Statement
\section{Problem Statement and Contribution}
%Different disciplines have varying naming conventions. In engineering, thesis tend to refer to problem(s) to be solved where other disciplines talk in terms of question(s) to be answered. In either case, this section has three main parts:

%\begin{itemize}
  %\item A concise statement of the question(s) that your thesis shall tackle.
  %\item Justification, by direct reference to literature review from previous
  %section that your question is previously unanswered.
 % \item Discussion of why it is worthwhile to answer this question.
%\end{itemize}

%Since this is one of the sections that the readers are definitely looking for, describe the issues more clearly by dividing this section into multiple subsections such as a spate section for ``Aims", another for listing ``Original Contributions", and yet another for mentioning the ``Limitations".


(Alvarez, 2020) "As can be inferred, the selection, implementation, operation and
maintenance of a System for Automatic Diagnosis of Failures is not a simple task, requiring at
each stage, care so that the result provided by the system, after its implementation, is within
the one initially specified.” Is it therefore possible to detect system faults within a storage
engine and rectify such faults to increase the quality and effectiveness of the storage engine,
by focusing on the reliability and correctness of the storage engine, given its complexity?
Through experience, the researcher has come to the understanding that all software design, algorithms and
coding processes contains one or more errors, some which are easy to detect while others not.

%Skipping Paragraphs by one line
\setlength{\parskip}{10pt}

If there was a method to avoid errors in the core system processes and development stages,
what would that look like? And how effective would it be? This dilemma lead Professor Leslie
Lamport to seek and develop an alternative approach to solving complex systems through
formal methods. “When I developed TLA, I realized that, for the first time, I had a formalism
that really was completely formal, so formal that mechanically checking TLA proofs should be
straightforward.” (Urban Engberg, Leslie Lamport, Peter Grønning, 1992).

%Skipping Paragraphs by one line
\setlength{\parskip}{10pt}

The purpose of this research is therefore to actively use TLA+ as a formal method in the
developing process of the storage engine to specify and evaluate the error reduction rate, by
focusing on the safety and correctness of the system of the storage engine. Through the use
of TLA+ the researcher aims to formally specify the critical components and or operations of
the system in order to verify its property against an implementation.

\section{RESEARCH OBJECTIVES}
The researcher aims to apply the TLA+ to the CRUD key operations of the storage engine. To
verify the properties of the implementation against the specification, during the development
and testing phase of the storage engine.
With the acronym CRUD the researcher aims to:
\begin{enumerate}

\item C = Create:
This will be the write action and subsequent write operations
\item R = Read:
This will be the read request sent by a client
\item U = Update:
This will be a write new version operation
\item D = Soft Delete:
This will be the stop write operation on a given version tree only.
\end{enumerate}
With reference to the above the researcher intends to:
\begin{enumerate}[(a)] % (a), (b), (c), ...
\item Specify the CRUD operations of the storage engine in TLA+.
\item Evaluate the design using the specification
\item Verify an existing implementation against the specification

\end{enumerate}




\chapter{Literature Review}\label{c-review}
Here you review the state of the art relevant to your thesis proposal. The idea is to present the major ideas in the state of the art right up to, but not including, your own personal brilliant ideas.

Critical analysis and comparisons should be made by pointing out the weakness of existing solutions and strengths of your proposal. You organize this section by idea, and not by author or by publication.

In certain situations, a background of the underlying concepts is required for better understanding of the research problem and also to improve the flow of the thesis. This could either be made an introductory part of this section or separately written in a prior section.





\begin{table}[!th]
\centering
% Use, for example, p{3.5cm} style for fixed sized columns
% consider using vbox to ensure large text is wrapped inside a column

%Reference list included

\begin{tabular}{|p{3cm}|p{9cm}|}
\hline
\textbf{Reference Type} & \textbf{Citation}\\ \hline
Article & \cite{Hayes1993}\\ \hline
Book & \cite[p.127-133]{Lamport_2005}\\ \hline
InProceedings & \cite*{Nicholls_1987}\\ \hline
InCollection & \cite{Lund_2019}\\ \hline
PhD Dissertation & \cite{Konnov_2019}\\ \hline
Masters Thesis & \cite{Lamport_2005}\\ \hline
Technical Report & \cite{Engberg1993}\\ \hline
Misc & \cite{Hayes1993}\\ \hline
\end{tabular}
\caption{Citation Styles.}
\label{t-References}
\end{table}


\chapter{Storage Engine}
A brief description of the storage engine an example use, replication processes and operations specific to this document

\section{Overview}
The Storage engine itself is not is not what this research paper is focused. However This paper relates to the storage engine by addressing some core issues in the development process to avoid potentially severe costs and termination due to one or more errors or conflicts at a later stage in the development process. 

\subsection{Functionality}
The researcher assumes that the storage engine should be able to;
\begin{itemize}
    \item Store and Retrieve data.
    \item Modify existing Data
    \item Keep all versions of the modification
    \item Delete the modification process
    \item Keep track of all modification
    \item Retrieve all or a certain version of the modification.
    \item Create several replica of the data
    \item implement a B+ Tree Data structure to handle the data
    \item Keep record of different data objects 
    \item Allow modification of on each object individually.
    \item Allow concurrent read write operations to occur. 
    \item Furthermore some other complex operations not yet defined in the scope of this paper. 
\end{itemize}

\section{Use Case 1}

\section{Use Case 2}

\section{Replication}

\section{Key Operations}
The key operation can be classified as the basic CRUD operations. 

\begin{itemize}
    \item Create: This is the first write operation.
    \item Read: This is all the reading operations.
    \item Update: This is the version modification and also a write operation
    \item Delete: This should not delete the data, but only prevent any further updates (versions and modifications) from being 
\end{itemize}





\chapter{Storage Engine Specification}

\section{Model Description}

\section{Constant and Variables}

\section{Initialization}

\section{Next State}

%\section{}


\section{Operation Specification }

\subsection{Write Operation}

\subsection{Read Operation}

\subsection{Update Operation}

\subsection{Delete Operation}

\label{c-methods}

This part of the thesis is much more free-form. It may have several subsections. But it all has only one purpose: to convince the examiners about the answer to the research question(s) or solution to the problem(s) that you set for yourself in the proposal.

\begin{definition}[Testing 1,2,3]
This definition is placed within a chapter so is its number.
\end{definition}

\begin{figure}[http]
\begin{center}
  \includegraphics[width=0.7\columnwidth]{nust.jpg}
  \caption{NUST Emblem.}
  \label{f-nust}
\end{center}
\end{figure}


\chapter{Verification of Implementation}
\label{c-results}

So show what you did so far (implementation and testing) that is relevant to
answering the question(s) or solving the problem(s).


%%!TEX root = ../thesis.tex
\chapter{Discussion} % (fold)
\label{chap:discussion}

\section{Section 1} % (fold)
\label{sec:section_11}
Led ut perspiciatis unde omnis iste natus error sit voluptatem accusantium doloremque laudantium, totam rem aperiam, eaque ipsa quae ab illo inventore veritatis et quasi architecto beatae vitae dicta sunt explicabo. Nemo enim ipsam voluptatem quia voluptas sit aspernatur aut odit aut fugit, sed quia consequuntur magni dolores eos qui ratione voluptatem sequi nesciunt. Neque porro quisquam est, qui dolorem ipsum quia dolor sit amet, consectetur, adipisci velit, sed quia non numquam eius modi tempora incidunt ut labore et dolore magnam aliquam quaerat voluptatem. Ut enim ad minima veniam, quis nostrum exercitationem ullam corporis suscipit laboriosam, nisi ut aliquid ex ea commodi consequatur? Quis autem vel eum iure reprehenderit qui in ea voluptate velit esse quam nihil molestiae consequatur, vel illum qui dolorem eum fugiat quo voluptas nulla pariatur?

\par
Led ut perspiciatis unde omnis iste natus error sit voluptatem accusantium doloremque laudantium, totam rem aperiam, eaque ipsa quae ab illo inventore veritatis et quasi architecto beatae vitae dicta sunt explicabo. Nemo enim ipsam voluptatem quia voluptas sit aspernatur aut odit aut fugit, sed quia consequuntur magni dolores eos qui ratione voluptatem sequi nesciunt. Neque porro quisquam est, qui dolorem ipsum quia dolor sit amet, consectetur, adipisci velit, sed quia non numquam eius modi tempora incidunt ut labore et dolore magnam aliquam quaerat voluptatem. Ut enim ad minima veniam, quis nostrum exercitationem ullam corporis suscipit laboriosam, nisi ut aliquid ex ea commodi consequatur? Quis autem vel eum iure reprehenderit qui in ea voluptate velit esse quam nihil molestiae consequatur, vel illum qui dolorem eum fugiat quo voluptas nulla pariatur?

\subsection{Subsection 1} % (fold)
\label{sub:subsection_111}

Led ut perspiciatis unde omnis iste natus error sit voluptatem accusantium doloremque laudantium, totam rem aperiam, eaque ipsa quae ab illo inventore veritatis et quasi architecto beatae vitae dicta sunt explicabo. Nemo enim ipsam voluptatem quia voluptas sit aspernatur aut odit aut fugit, sed quia consequuntur magni dolores eos qui ratione voluptatem sequi nesciunt. Neque porro quisquam est, qui dolorem ipsum quia dolor sit amet, consectetur, adipisci velit, sed quia non numquam eius modi tempora incidunt ut labore et dolore magnam aliquam quaerat voluptatem (see Fig. \ref{fig:demo_figure_1}). Ut enim ad minima veniam, quis nostrum exercitationem ullam corporis suscipit laboriosam, nisi ut aliquid ex ea commodi consequatur? Quis autem vel eum iure reprehenderit qui in ea voluptate velit esse quam nihil molestiae consequatur, vel illum qui dolorem eum fugiat quo voluptas nulla pariatur? \cite{bruno,smith}

\begin{figure}[h]
\centering
\includegraphics[width=0.5\textwidth]{./img/350x150.png}
\caption{The caption of the figure goes here.}
\label{fig:demo_figure_1}
\end{figure}

\par
Lorem ipsum dolor sit amet, consectetur adipiscing elit, sed do eiusmod tempor incididunt ut labore et dolore magna aliqua. Ut enim ad minim veniam, quis nostrud exercitation ullamco laboris nisi ut aliquip ex ea commodo consequat \cite{bruno}. Duis aute irure dolor in reprehenderit in voluptate velit esse cillum dolore eu fugiat nulla pariatur. (see Fig. \ref{fig:demo_figure_2}) Excepteur sint occaecat cupidatat non proident, sunt in culpa qui officia deserunt mollit anim id est laborum. \cite{smith}
% subsection subsection_111 (end)

\subsection{Subsection 2} % (fold)
\label{sub:subsection_112}
Quis autem vel eum iure reprehenderit qui in ea voluptate velit esse quam nihil molestiae consequatur, vel illum qui dolorem eum fugiat quo voluptas nulla pariatur. Ut enim ad minima veniam, quis nostrum exercitationem ullam corporis suscipit laboriosam, nisi ut aliquid ex ea commodi consequatur? Quis autem vel eum iure reprehenderit qui in ea voluptate velit esse quam nihil molestiae consequatur, vel illum qui dolorem eum fugiat quo voluptas nulla pariatur.

\begin{figure}[h]
\centering
\includegraphics[width=0.5\textwidth]{./img/500X300.png}
\caption{The caption of the figure goes here.}
\label{fig:demo_figure_2}
\end{figure}

\par
Quis autem vel eum iure reprehenderit qui in ea voluptate velit esse quam nihil molestiae consequatur, vel illum qui dolorem eum fugiat quo voluptas nulla pariatur. Ut enim ad minima veniam, quis nostrum exercitationem ullam corporis suscipit laboriosam, nisi ut aliquid ex ea commodi consequatur? Quis autem vel eum iure reprehenderit qui in ea voluptate velit esse quam nihil molestiae consequatur, vel illum qui dolorem eum fugiat quo voluptas nulla pariatur.
% subsection subsection_112 (end)
% section section_11 (end)

\section{Section 2} % (fold)
\label{sec:section_12}
Quis autem vel eum iure reprehenderit qui in ea voluptate velit esse quam nihil molestiae consequatur, vel illum qui dolorem eum fugiat quo voluptas nulla pariatur. Ut enim ad minima veniam, quis nostrum exercitationem ullam corporis suscipit laboriosam, nisi ut aliquid ex ea commodi consequatur? Quis autem vel eum iure reprehenderit qui in ea voluptate velit esse quam nihil molestiae consequatur, vel illum qui dolorem eum fugiat quo voluptas nulla pariatur.

\par
Quis autem vel eum iure reprehenderit qui in ea voluptate velit esse quam nihil molestiae consequatur, vel illum qui dolorem eum fugiat quo voluptas nulla pariatur. Ut enim ad minima veniam, quis nostrum exercitationem ullam corporis suscipit laboriosam, nisi ut aliquid ex ea commodi consequatur? Quis autem vel eum iure reprehenderit qui in ea voluptate velit esse quam nihil molestiae consequatur, vel illum qui dolorem eum fugiat quo voluptas nulla pariatur.

\subsection{Subsection 1} % (fold)
\label{sub:subsection_121}
Quis autem vel eum iure reprehenderit qui in ea voluptate velit esse quam nihil molestiae consequatur, vel illum qui dolorem eum fugiat quo voluptas nulla pariatur. Ut enim ad minima veniam, quis nostrum exercitationem ullam corporis suscipit laboriosam, nisi ut aliquid ex ea commodi consequatur? Quis autem vel eum iure reprehenderit qui in ea voluptate velit esse quam nihil molestiae consequatur, vel illum qui dolorem eum fugiat quo voluptas nulla pariatur.

\par
Quis autem vel eum iure reprehenderit qui in ea voluptate velit esse quam nihil molestiae consequatur, vel illum qui dolorem eum fugiat quo voluptas nulla pariatur. Ut enim ad minima veniam, quis nostrum exercitationem ullam corporis suscipit laboriosam, nisi ut aliquid ex ea commodi consequatur? Quis autem vel eum iure reprehenderit qui in ea voluptate velit esse quam nihil molestiae consequatur, vel illum qui dolorem eum fugiat quo voluptas nulla pariatur.
% subsection subsection_121 (end)

\subsection{Subsection 2} % (fold)
\label{sub:subsection_122}
Quis autem vel eum iure reprehenderit qui in ea voluptate velit esse quam nihil molestiae consequatur, vel illum qui dolorem eum fugiat quo voluptas nulla pariatur. Ut enim ad minima veniam, quis nostrum exercitationem ullam corporis suscipit laboriosam, nisi ut aliquid ex ea commodi consequatur? Quis autem vel eum iure reprehenderit qui in ea voluptate velit esse quam nihil molestiae consequatur, vel illum qui dolorem eum fugiat quo voluptas nulla pariatur.

\par
Quis autem vel eum iure reprehenderit qui in ea voluptate velit esse quam nihil molestiae consequatur, vel illum qui dolorem eum fugiat quo voluptas nulla pariatur. Ut enim ad minima veniam, quis nostrum exercitationem ullam corporis suscipit laboriosam, nisi ut aliquid ex ea commodi consequatur? Quis autem vel eum iure reprehenderit qui in ea voluptate velit esse quam nihil molestiae consequatur, vel illum qui dolorem eum fugiat quo voluptas nulla pariatur.
% subsection subsection_122 (end)
% section section_12 (end)
%!TEX root = ../thesis.tex
\chapter{Conclusion} % (fold)
\label{cha:_literature_review}

\section{Section 1} % (fold)
\label{sec:section_21}
Led ut perspiciatis unde omnis iste natus error sit voluptatem accusantium doloremque laudantium, totam rem aperiam, eaque ipsa quae ab illo inventore veritatis et quasi architecto beatae vitae dicta sunt explicabo. Nemo enim ipsam voluptatem quia voluptas sit aspernatur aut odit aut fugit, sed quia consequuntur magni dolores eos qui ratione voluptatem sequi nesciunt. Neque porro quisquam est, qui dolorem ipsum quia dolor sit amet, consectetur.

\par
Led ut perspiciatis unde omnis iste natus error sit voluptatem accusantium doloremque laudantium, totam rem aperiam, eaque ipsa quae ab illo inventore veritatis et quasi architecto beatae vitae dicta sunt explicabo. Nemo enim ipsam voluptatem quia voluptas sit aspernatur aut odit aut fugit, sed quia consequuntur magni dolores eos qui ratione voluptatem sequi nesciunt. Neque porro quisquam est, qui dolorem ipsum quia dolor sit amet, consectetur.

\subsection{Subsection 1} % (fold)
\label{sub:subsection_211}
Led ut perspiciatis unde omnis iste natus error sit voluptatem accusantium doloremque laudantium, totam rem aperiam, eaque ipsa quae ab illo inventore veritatis et quasi architecto beatae vitae dicta sunt explicabo. Nemo enim ipsam voluptatem quia voluptas sit aspernatur aut odit aut fugit, sed quia consequuntur magni dolores eos qui ratione voluptatem sequi nesciunt. Neque porro quisquam est, qui dolorem ipsum quia dolor sit amet, consectetur.

\begin{equation} \label{eq1}
  \begin{split}
    A & = \frac{\pi r^2}{2} \\
      & = \frac{1}{2} \pi r^2
  \end{split}
\end{equation}

\par
Adipisci velit, sed quia non numquam eius modi tempora incidunt ut labore et dolore magnam aliquam quaerat voluptatem. Ut enim ad minima veniam, quis nostrum exercitationem ullam corporis suscipit laboriosam, nisi ut aliquid ex ea commodi consequatur? Quis autem vel eum iure reprehenderit qui in ea voluptate velit esse quam nihil molestiae consequatur, vel illum qui dolorem eum fugiat quo voluptas nulla pariatur? \cite{bruno,smith}

\begin{table}[!htb]
\centering
    \begin{tabular}{xyyy}
    \hline
    \rowcolor{Gray}
    \textsc{Column 1} & \textsc{Column 2} & \textsc{Column 3} & \textsc{Column 4}\\
    \hline
    1, 1 & 1, 2 & 1, 3 & 1, 4 \\
    2, 1 & 2, 2 & 2, 3 & 2, 4 \\
    3, 1 & 3, 2 & 3, 3 & 3, 4 \\
    4, 1 & 4, 2 & 4, 3 & 4, 4 \\ \hline
    \end{tabular}
    \caption{The caption of the table goes here.}
    \label{tab:demo_table}
\end{table}

\par
Lorem ipsum dolor sit amet, consectetur adipiscing elit, sed do eiusmod tempor incididunt ut labore et dolore magna aliqua. Ut enim ad minim veniam, quis nostrud exercitation ullamco\footnote{\label{foot_note_1}Led ut perspiciatis unde omnis iste natus error sit voluptatem accusantium doloremque laudantium, totam rem aperiam.} laboris nisi ut aliquip ex ea commodo consequat \cite{bruno}. Duis aute irure dolor in reprehenderit in voluptate velit esse cillum dolore eu fugiat nulla pariatur. Excepteur sint occaecat cupidatat non proident, sunt in culpa qui officia deserunt mollit anim id est laborum. \cite{smith}
\begin{multline}
p(x) = 3x^6 + 14x^5y + 590x^4y^2 + 19x^3y^3\\
- 12x^2y^4 - 12xy^5 + 2y^6 - a^3b^3
\end{multline}
% subsection subsection_211 (end)

\subsection{Subsection 2} % (fold)
\label{sub:subsection_212}
Led ut perspiciatis unde omnis iste natus error sit voluptatem accusantium doloremque laudantium, totam rem aperiam, eaque ipsa quae ab illo inventore veritatis et quasi architecto beatae vitae dicta sunt explicabo. Nemo enim ipsam voluptatem quia voluptas sit aspernatur aut odit aut fugit, sed quia consequuntur magni dolores eos qui ratione voluptatem sequi nesciunt. Neque porro quisquam est, qui dolorem ipsum quia dolor sit amet, consectetur.

\par
Led ut perspiciatis unde omnis iste natus error sit voluptatem accusantium doloremque laudantium, totam rem aperiam, eaque ipsa quae ab illo inventore veritatis et quasi architecto beatae vitae dicta sunt explicabo. Nemo enim ipsam voluptatem quia voluptas sit aspernatur aut odit aut fugit, sed quia consequuntur magni dolores eos qui ratione voluptatem sequi nesciunt. Neque porro quisquam est, qui dolorem ipsum quia dolor sit amet, consectetur.
% subsection subsection_212 (end)
% section section_21 (end)

\section{Section 2} % (fold)
\label{sec:section_22}
Led ut perspiciatis unde omnis iste natus error sit voluptatem accusantium doloremque laudantium, totam rem aperiam, eaque ipsa quae ab illo inventore veritatis et quasi architecto beatae vitae dicta sunt explicabo. Nemo enim ipsam voluptatem quia voluptas sit aspernatur aut odit aut fugit, sed quia consequuntur magni dolores eos qui ratione voluptatem sequi nesciunt. Neque porro quisquam est, qui dolorem ipsum quia dolor sit amet, consectetur.

\par
Led ut perspiciatis unde omnis iste natus error sit voluptatem accusantium doloremque laudantium, totam rem aperiam, eaque ipsa quae ab illo inventore veritatis et quasi architecto beatae vitae dicta sunt explicabo. Nemo enim ipsam voluptatem quia voluptas sit aspernatur aut odit aut fugit, sed quia consequuntur magni dolores eos qui ratione voluptatem sequi nesciunt. Neque porro quisquam est, qui dolorem ipsum quia dolor sit amet, consectetur.

\section{Future Work} % (fold)
\label{sub:subsection_221}
Led ut perspiciatis unde omnis iste natus error sit voluptatem accusantium doloremque laudantium, totam rem aperiam, eaque ipsa quae ab illo inventore veritatis et quasi architecto beatae vitae dicta sunt explicabo. Nemo enim ipsam voluptatem quia voluptas sit aspernatur aut odit aut fugit, sed quia consequuntur magni dolores eos qui ratione voluptatem sequi nesciunt. Neque porro quisquam est, qui dolorem ipsum quia dolor sit amet, consectetur.

\par
Led ut perspiciatis unde omnis iste natus error sit voluptatem accusantium doloremque laudantium, totam rem aperiam, eaque ipsa quae ab illo inventore veritatis et quasi architecto beatae vitae dicta sunt explicabo. Nemo enim ipsam voluptatem quia voluptas sit aspernatur aut odit aut fugit, sed quia consequuntur magni dolores eos qui ratione voluptatem sequi nesciunt. Neque porro quisquam est, qui dolorem ipsum quia dolor sit amet, consectetur.
% subsection subsection_221 (end)

%\bibliography{references}
%\bibliographystyle{apa}
%\nocite{*}
\printbibliography


%\begin{comment}

--------------------------- MODULE StorageObject ---------------------------
(*
Fot the Object Specification We will use the Spaning Tree Algorithm and modify it to suit our needs 
The Spanning tree Allows for the Groth and also control and identifying neighbors 

This can be changed to parent-chlid and will allow operations shuch as 
1. update (promotion to parent) and create child
2. Delete (remove child and downgrade from parent to child and change to done)
3. write ( will only be done on initial root as all children must come from root node)
4. read (locate the child ( Reference parents and Grand parents { incremental update read }

The naming conventions are subject to change if need be 

The object will have two control states
controState and subControlState

controState: This will be a set of initial States ( Start, middle and done)
    Start: indicates that no current version exists and must be created
    middle: indicates that child nodes (versions ) exists and can be updated
    done: indicated that:
        1. No more child nodes can be created
        2. Terminationi has occured 
        3. Hopefully return last version as Child of 

subControlState:
    Lf: indicates a leaf node (last child) {last version on that branch}
    Pt: indicates a node became a parent
    Co: indicates  back tracking neighbor nodes to the root. ( list of app parents conected to current node)


*)

EXTENDS Integers, FiniteSets, Randomization                      \*introduce Integers modules and Naturals*\

CONSTANTS Nodes, Root, MaxCardinality    \* we use MaxCardinality instead of Infinity*\

(*
for parent nodes including the root we will denote as key
for edges we will denote as ver (version)

node n has a neighbor m such that ver(n) to ver(m) + 1 or ver(m) + x where x is an positive integer 
*)

VARIABLE object, controlState, subControlState, key, Ver, Edges, Path

TypeOK == /\ controlState \in {"Start", "middle", "done"}
          /\ subControlState \in {"Lf", "Pt", "Co"}
          /\ key  \in [Nodes -> Nodes]
          /\ Ver \in [Nodes -> Nat]
          /\ Path \in [Paths -> Paths]    
          
ASSUME /\ Root \in Nodes
       /\ MaxCardinality \in Nat
       /\ MaxCardinality >= Cardinality(Nodes)

(* We will denote versions as a set of neighbors of node n in the graph joined by an edge to n
*)

Versions(n) == {m \in Nodes : {m, n} \in Edges}


(*
wE NOW SPECIFY THE ALGORITHM 
*) 

vars == <<key, Ver, Edges>> \* or child nodes instead Childnodes (how many child nodes is linked to  current parent *\

         
\* initial write function*\

Init == /\ key = [n \in Nodes |-> n]
        /\ Ver = [n \in Nodes |-> IF n = Root THEN 0 ELSE MaxCardinality]
        /\ Edges \in {E \in SUBSET (SUBSET Nodes) : \A e \in E : Cardinality(e) = 2}
        /\ controlState = "middle"
        /\ subControlState = "Pt" 

(* All Methods to incorporate

vars == <<timestamp, values, deliverQueues>>
nodeIds == 1..N_NODES

\Prepares to read   *\
DeliverSet(n, T, t, k, v) ==
  values' = [
      values EXCEPT ![n] = {<<tp, kp, vp>> \in values[n] : tp \notin T} \union { <<t, k, v>> }
    ]

\* Clears current read req *\
DeliverDelete(n, T) ==
  values' = [
    values EXCEPT ![n] = {<<t, k, v>> \in values[n] : t \notin T}
  ]

\* Checks read Commad and Parents connected to it *\
Deliver(n, command, payload) ==
  \/ command = "set"
    /\ DeliverSet(n, payload[1], payload[2], payload[3], payload[4])
  \/ command = "delete"
    /\ DeliverDelete(n, payload)

\* Prints Tree Graph with Children *\
Broadcast(n, command, payload) ==
  /\ Deliver(n, command, payload)
  /\ deliverQueues' = [
      i \in nodeIds |->
        IF i = n THEN
          deliverQueues[i]
        ELSE
          Append(deliverQueues[i], <<command, payload>>)
      ]

\* UpdateWrite Checks if not deleted and if can be updated *\
RequestSet(n, k, v) ==  
  LET matches == {<<t, kp, vp>> \in values[n] : k = kp}  IN
  LET T == {t : <<t, kp, vp>> \in matches}  IN
    /\ timestamp' = timestamp + 1
    /\ Broadcast(n, "set", <<T, timestamp, k, v>>)

\* Checks if we can Delete *\
RequestDelete(n, k) ==
  LET matches == {<<t, kp, v>> \in values[n] : k = kp}  IN
  LET T == {t : <<t, kp, v>> \in matches}  IN
    /\ T /= {}
    /\ Broadcast(n, "delete", T)


\* UpdateWrite Action *\
RequestSetOnNode ==
  /\ timestamp < MAX_TIMESTAMP
  /\ \E <<n, k, v>> \in nodeIds \X KEYS \X VALUES : RequestSet(n, k, v)

\* Checks vurrent Version and if node exists *\
RequestDeleteOnNode ==
  /\ \E <<n, k>> \in nodeIds \X KEYS : RequestDelete(n, k)
  /\ UNCHANGED timestamp


\* ReadNote *\
DeliverOnNode ==
  \E n \in nodeIds :
    /\ Len(deliverQueues[n]) > 0
    /\ \E <<command, payload>> \in {Head(deliverQueues[n])} :
        Deliver(n, command, payload)
    /\ deliverQueues' = [deliverQueues EXCEPT ![n] = Tail(deliverQueues[n])]
  /\ UNCHANGED timestamp

\* Checks if there are not Children nodes *\
DeliverQueuesIsEmpty ==
  \A n \in nodeIds: Len(deliverQueues[n]) = 0


\* Deleting Action  *\
Terminating ==
  /\ DeliverQueuesIsEmpty
  /\ UNCHANGED vars 


*)

        
(* I need to define my next state based on the functions I will define
    Write
    Update
    Delete
    Read
    
    as follows
    
    Next == \/ writeRoot
            \/ UpdateWrte
            \/ ReadNode
            \/ DeleteNdoe
            
    /* my spec will then be *\
    
    Spec == Init /\ [][Next]_var /\Wf_vars(ReadNode)



*)        
Next == \E n \in Nodes :
          \E m \in Versions(n) : 
             /\ Ver[m] < 1 + Ver[n]
             /\ \E d \in (Ver[m]+1) .. (Ver[n] - 1) :
                    /\ Ver' = [Ver EXCEPT ![n] = d]
                    /\ key'  = [key  EXCEPT ![n] = m]


\* Inital call*\
Spec == Init /\ [][Next]_vars /\ WF_vars(Next)
(*
\* update Function*\

UpdateWrite(key, ver) == Init /\ [][Next]_vars /\ WF_vars(Next)

\*Delete Function *\


\*Update unction *\

*)

\*Checking if the parent nodes are correct *\
PostCondition == 
  \A n \in Nodes :
    \/ /\ n = Root 
       /\ Ver[n] = 0
       /\ key[n] = n
    \/ /\ Ver[n] = MaxCardinality 
       /\ key[n] = n
       /\ \A m \in Versions(n) : Ver[m] = MaxCardinality
    \/ /\ Ver[n] \in 1..(MaxCardinality-1)
       /\ key[n] \in Versions(n)
       /\ Ver[n] = Ver[key[n]] + 1






Safety == []((~ ENABLED Next) => PostCondition)




Liveness == <>(~ ENABLED Next) 
























 
=============================================================================
\* Modification History
\* Last modified Tue Oct 25 11:00:17 CAST 2022 by goodwill
\* Created Tue Oct 04 09:12:06 CAST 2022 by goodwill

\end{comment}
\begin{appendix}
\chapter{First Appendix}
The separate numbering of appendices is also supported by LaTeX. The \textit{appendix} macro can be used to indicate that following chapters are to be numbered as appendices. Only use the \textit{appendix} macro once for all appendices.
\end{appendix}

\pagenumbering{gobble}
\begin{table}[]
\begin{tabular}{lll}
               & \multicolumn{1}{r}{Annex 'A'}                                                                                   &  \\
               & \multicolumn{1}{r}{office order: 0986/29/ACB/SEECS}                                                             &  \\
               & \multicolumn{1}{r}{Date \_\_\_\_\_\_\_Oct, 2009}                                                                &  \\
\multicolumn{2}{l}{\textbf{Th.ECL (MS Thesis   Evaluation Check List)}}                                                          &  \\
               &                                                                                                                 &  \\
\multicolumn{2}{l}{Student  Name:}                                                                                               &  \\
\multicolumn{2}{l}{Registration:}                                                                                                &  \\
               &                                                                                                                 &  \\
\multicolumn{2}{l}{\textbf{Cover and title   page of the thesis}}                                                                &  \\
T1.            & Student's name and registration   number is written.                                                            &  \\
T2.            & Supervisor's name is   mentioned.                                                                               &  \\
T3.            & Title of the degree   is written correctly.                                                                     &  \\
T4.            & University and   school's name are written correctly.                                                           &  \\
T5.            & Date of   completion/defense (only year and month) is mentioned.                                                &  \\
               &                                                                                                                 &  \\
\multicolumn{2}{l}{\textbf{Style and formatting issues}}                                                                         &  \\
S1.            & Consistent font (Times New Roman) is   used throughout the thesis.                                              &  \\
S2.            & Page numbering is   done appropriately.                                                                         &  \\
S3.            & Figures are readable   and are aligned correctly.                                                               &  \\
S4.            & Captions for tables   and figures use consistent format and style.                                              &  \\
S5.            & Table of   Contents/Figures/Tables follow proper indentation/styling.                                           &  \\
S6.            & Chapter name and   numbering follows consistent style.                                                          &  \\
               &                                                                                                                 &  \\
\multicolumn{2}{l}{\textbf{References/Bibliography}}                                                                             &  \\
R1.            & References are   sorted on last name of authors (or in the order of citation in the text).                      &  \\
R2.            & References follow   consistent style such as ACM or IEEE-Tran.                                                  &  \\
R3.            & Mandatory slots of   references are filled correctly (such as Author, Title, Journal, Year).                    &  \\
               &                                                                                                                 &  \\
\multicolumn{2}{l}{\textbf{General Issues}}                                                                                      &  \\
G1.            & Certificate of Originality signed by   the student is present.                                                  &  \\
G2.            & Plagiarism   report (from Euphorus) signed by supervisor is presented along with the   thesis.                  &  \\
G3.            & Thesis is submitted   within allowed time span for completion of thesis.                                        &  \\
               &                                                                                                                 &  \\
\multicolumn{2}{l}{\textbf{Abstract (Note:   This section covers only the abstract of the thesis)}}                              &  \\
A1.            & There are no typing or grammatic   mistakes in the abstract.                                                    &  \\
A2.            & Problem statement is   clearly mentioned.                                                                       &  \\
A3.            & Background to problem   statement is also explained.                                                            &  \\
A4.            & Startling statement   (preferably a paragraph) about the thesis/hypothesis is present.                          &  \\
A5.            & Implication of the   startling statement is demonstrated briefly.                                               &  \\
               &                                                                                                                 &  \\
\multicolumn{2}{l}{\textbf{Results, Evaluation, and Conclusion}}                                                                 &  \\
E1.            & Research is   validated either empirically or analytically (Note: This doesn’t cover   quality of the results). &  \\
E2.            & Outcome of this   thesis is contrasted with other similar research initiatives.                                 &  \\
E3.            & Significance of this   research is discussed in appropriate length.                                             &  \\
\multicolumn{2}{l}{\textbf{Thesis Format}}                                                                                       &  \\
\textbf{Sno}   & \multicolumn{2}{l}{\textbf{HQ NUST Format}}                                                                        \\
1              & Title Page                                                                                                      &  \\
2              & Thesis Acceptance   Certificate                                                                                 &  \\
3              & Approval Page                                                                                                   &  \\
4              & Dedicatoin                                                                                                      &  \\
5              & Certificate of   Originality                                                                                    &  \\
6              & Acknowledgement                                                                                                 &  \\
7              & Table of Contents                                                                                               &  \\
8              & List of Abbreviation                                                                                            &  \\
9              & List of Tables                                                                                                  &  \\
10             & List of Figures                                                                                                 &  \\
11             & Abstract                                                                                                        &  \\
12             & Main Body                                                                                                       &  \\
\multicolumn{2}{l}{\textbf{Checklist for Components in Main Body}}                                                               &  \\
Sno            & HQ NUST Fromat                                                                                                  &  \\
1              & Introduction                                                                                                    &  \\
2              & Literature Review                                                                                               &  \\
3              & Methodology                                                                                                     &  \\
4              & Results                                                                                                         &  \\
5              & Discussion                                                                                                      &  \\
6              & Conclusion                                                                                                      &  \\
7              & Recommandation                                                                                                  &  \\
8              & Reference                                                                                                       &  \\
9              & Appendices                                                                                                      &  \\
10             & Index (Optional)                                                                                                &  \\
               &                                                                                                                 &  \\
\multicolumn{2}{l}{\textbf{Additional Remarks:}}                                                                                 &  \\ \cline{1-2}
\multicolumn{2}{|l|}{}                                                                                                           &  \\ \cline{1-2}
\multicolumn{2}{|l|}{}                                                                                                           &  \\ \cline{1-2}
\multicolumn{2}{|l|}{}                                                                                                           &  \\ \cline{1-2}
\multicolumn{2}{l}{\textbf{OiC   MS Thesis:}}                                                                                    &  \\
\textbf{Date:} &                                                                                                                 & 
\end{tabular}
\end{table}
\begin{table}[!ht]
\footnotesize
\begin{adjustbox}{width=\linewidth,center}
\begin{tabular}{p{0.6cm}p{9cm}p{0.6cm}}
\multicolumn{3}{r}{Annex 'A'}                                                                                                                                                            \\
\multicolumn{3}{r}{office order: 0986/29/ACB/SEECS}                                                                                                                                      \\
\multicolumn{3}{r}{Date \_\_\_\_,\_\_\_,20\_\_\_}                                                                                                                                        \\
\multicolumn{2}{c}{\textbf{Th.ECL (MS Thesis   Evaluation Check List)}}                                                                                          &                       \\
                          &                                                                                                                                      &                       \\ \midrule
\multicolumn{3}{|l|}{Student  Name:}                                                                                                                                                     \\ \midrule
\multicolumn{3}{|l|}{Registration:}                                                                                                                                                      \\ \midrule
                          &                                                                                                                                      &                       \\ \midrule
\multicolumn{3}{|l|}{\textbf{Cover and title   page of the thesis}}                                                                                                                      \\ \midrule
\multicolumn{1}{|l|}{T1.} & \multicolumn{1}{l|}{Student's name and registration   number is written.}                                                            & \multicolumn{1}{l|}{} \\ \midrule
\multicolumn{1}{|l|}{T2.} & \multicolumn{1}{l|}{Supervisor's name is   mentioned.}                                                                               & \multicolumn{1}{l|}{} \\ \midrule
\multicolumn{1}{|l|}{T3.} & \multicolumn{1}{l|}{Title of the degree   is written correctly.}                                                                     & \multicolumn{1}{l|}{} \\ \midrule
\multicolumn{1}{|l|}{T4.} & \multicolumn{1}{l|}{University and   school's name are written correctly.}                                                           & \multicolumn{1}{l|}{} \\ \midrule
\multicolumn{1}{|l|}{T5.} & \multicolumn{1}{l|}{Date of   completion/defense (only year and month) is mentioned.}                                                & \multicolumn{1}{l|}{} \\ \midrule
                          &                                                                                                                                      &                       \\ \midrule
\multicolumn{3}{|l|}{\textbf{Style and formatting issues}}                                                                                                                               \\ \midrule
\multicolumn{1}{|l|}{S1.} & \multicolumn{1}{l|}{Consistent font (Times New Roman) is   used throughout the thesis.}                                              & \multicolumn{1}{l|}{} \\ \midrule
\multicolumn{1}{|l|}{S2.} & \multicolumn{1}{l|}{Page numbering is   done appropriately.}                                                                         & \multicolumn{1}{l|}{} \\ \midrule
\multicolumn{1}{|l|}{S3.} & \multicolumn{1}{l|}{Figures are readable   and are aligned correctly.}                                                               & \multicolumn{1}{l|}{} \\ \midrule
\multicolumn{1}{|l|}{S4.} & \multicolumn{1}{l|}{Captions for tables   and figures use consistent format and style.}                                              & \multicolumn{1}{l|}{} \\ \midrule
\multicolumn{1}{|l|}{S5.} & \multicolumn{1}{l|}{Table of   Contents/Figures/Tables follow proper indentation/styling.}                                           & \multicolumn{1}{l|}{} \\ \midrule
\multicolumn{1}{|l|}{S6.} & \multicolumn{1}{l|}{Chapter name and   numbering follows consistent style.}                                                          & \multicolumn{1}{l|}{} \\ \midrule
                          &                                                                                                                                      &                       \\ \midrule
\multicolumn{3}{|l|}{\textbf{References/Bibliography}}                                                                                                                                   \\ \midrule
\multicolumn{1}{|l|}{R1.} & \multicolumn{1}{l|}{References are   sorted on last name of authors (or in the order of citation in the text).}                      & \multicolumn{1}{l|}{} \\ \midrule
\multicolumn{1}{|l|}{R2.} & \multicolumn{1}{l|}{References follow   consistent style such as ACM or IEEE-Tran.}                                                  & \multicolumn{1}{l|}{} \\ \midrule
\multicolumn{1}{|l|}{R3.} & \multicolumn{1}{l|}{Mandatory slots of   references are filled correctly (such as Author, Title, Journal, Year).}                    & \multicolumn{1}{l|}{} \\ \midrule
                          &                                                                                                                                      &                       \\ \midrule
\multicolumn{3}{|l|}{\textbf{General Issues}}                                                                                                                                            \\ \midrule
\multicolumn{1}{|l|}{G1.} & \multicolumn{1}{l|}{Certificate of Originality signed by   the student is present.}                                                  & \multicolumn{1}{l|}{} \\ \midrule
\multicolumn{1}{|l|}{G2.} & \multicolumn{1}{l|}{Plagiarism   report (from Euphorus) signed by supervisor is presented along with the   thesis.}                  & \multicolumn{1}{l|}{} \\ \midrule
\multicolumn{1}{|l|}{G3.} & \multicolumn{1}{l|}{Thesis is submitted   within allowed time span for completion of thesis.}                                        & \multicolumn{1}{l|}{} \\ \midrule
                          &                                                                                                                                      &                       \\ \midrule
\multicolumn{3}{|l|}{\textbf{Abstract (Note:   This section covers only the abstract of the thesis)}}                                                                                    \\ \midrule
\multicolumn{1}{|l|}{A1.} & \multicolumn{1}{l|}{There are no typing or grammatic   mistakes in the abstract.}                                                    & \multicolumn{1}{l|}{} \\ \midrule
\multicolumn{1}{|l|}{A2.} & \multicolumn{1}{l|}{Problem statement is   clearly mentioned.}                                                                       & \multicolumn{1}{l|}{} \\ \midrule
\multicolumn{1}{|l|}{A3.} & \multicolumn{1}{l|}{Background to problem   statement is also explained.}                                                            & \multicolumn{1}{l|}{} \\ \midrule
\multicolumn{1}{|l|}{A4.} & \multicolumn{1}{l|}{Startling statement   (preferably a paragraph) about the thesis/hypothesis is present.}                          & \multicolumn{1}{l|}{} \\ \midrule
\multicolumn{1}{|l|}{A5.} & \multicolumn{1}{l|}{Implication of the   startling statement is demonstrated briefly.}                                               & \multicolumn{1}{l|}{} \\ \midrule
                          &                                                                                                                                      &                       \\ \midrule
\multicolumn{3}{|l|}{\textbf{Results, Evaluation, and Conclusion}}                                                                                                                       \\ \midrule
\multicolumn{1}{|l|}{E1.} & \multicolumn{1}{l|}{Research is   validated either empirically or analytically (Note: This doesn’t cover   quality of the results).} & \multicolumn{1}{l|}{} \\ \midrule
\multicolumn{1}{|l|}{E2.} & \multicolumn{1}{l|}{Outcome of this   thesis is contrasted with other similar research initiatives.}                                 & \multicolumn{1}{l|}{} \\ \midrule
\multicolumn{1}{|l|}{E3.} & \multicolumn{1}{l|}{Significance of this   research is discussed in appropriate length.} \\ \midrule
\end{tabular}
\end{adjustbox}
\end{table}



\end{document}